% !TEX root = ../main.tex
%-----------------------------------------------------------------------
%  Appendix: verification before production.
%  \input by ../main.tex -- not compiled on its own.
%-----------------------------------------------------------------------
%-----------------------------------------------------------------------
\section{Verification before production}
%-----------------------------------------------------------------------

\subsection{Verifying a build}
\label{app:verify}

Applies to any new machine, and to this one if the build is ever regenerated.

\begin{enumerate}[nosep]
  \item Re-run the reference adjoint configuration (unperturbed $\kv$,
        \code{nIter0 = 3162240}) for a few hundred timesteps.
  \item Compare the cost function value and a few \code{ADJ*} fields against the
        reference run's output on scratch (job 28486).
  \item Expect agreement to near round-off. Large differences indicate a build or
        input problem rather than physics.
\end{enumerate}

A cheaper and stronger check exists for the forward model: a 10-year run from
rest with \code{from\_rest\_visc2x} reproduces the first 10 years of the 200-year
spin-up \textbf{bit-identically} --- 161 \code{dynDiag} field comparisons, all of
\code{surfDiag}, \code{atmDiag} and \code{viscDiag}, and 1334 monitor values
across 134 variables, every one at exactly zero difference. This should be run
first on a new build, as it catches a compiler or source change that alters the
physics.

\subsection{The gradient check cannot be used as a verification}
\label{app:grdchk}

\code{useGrdchk = .TRUE.} is set, so the check does run on every adjoint job,
its settings held in \code{input\_tap/data.grdchk}. It cannot pass at its
current setting, and it never has.

The check perturbs \code{xx\_theta} at
\code{iGloPos=4, jGloPos=8, kGloPos=1}, near the southern boundary, while the
cost section is at $j=127$. Sensitivity at that point is about
$6\times10^{-10}$ against a field maximum of about $3.9\times10^{-2}$, so the
finite difference measures run-to-run noise rather than the perturbation. Both
perturbed runs land on the \emph{same side} of the base cost, which a genuine
first derivative cannot produce. The check fails by roughly eight orders of
magnitude, on this build and on the May 2026 production run alike.

The cause is the location of the check, not the adjoint. To make it meaningful,
move the point into the sensitive region ---
\code{iGloPos=2, jGloPos=127, kGloPos=26} --- and raise \code{grdchk\_eps} to
about \code{1e-3} so the response clears the noise floor. That is worth doing
once. For the ensemble itself, the finite-difference validation in
Section~\ref{sec:fd} is the real check and a considerably stronger one.

\subsection{Benchmarking}

This machine is already measured; the figures are in
Section~\ref{sec:compute}. On a new machine, repeat the three measurements
there, then time four runs packed on one node ($4\times27 = 108$ of 128 cores)
against one run per node. A 27-rank job is billed for the whole node either way,
so if contention is mild, packing cuts charged node-hours by up to 4$\times$.
This does not apply here, where a run occupies its node.

